\VerdierPhysicalPageAnchor{0070}%
\VerdierSemanticAnchor{ch1-2-2-3-continuation}%
dénombrables sont représentables :

\VerdierSemanticAnchor{ch1-2-2-3-eq3}
\begin{equation}
  \begin{aligned}
    \left(\int_{\varphi}^{\Pi}Y\right)^{\tau}
      &=\prod_{\varphi(\sigma)=\tau}Y^{\sigma},\\
    d_{\int_{\varphi}^{\Pi}Y}^{\tau,i}
      &\quad\text{défini par la matrice }M(\tau,i,Y),\\[2pt]
    \left(\int_{\varphi}^{\Sigma}Y\right)^{\tau}
      &=\bigoplus_{\varphi(\sigma)=\tau}Y^{\sigma},\\
    d_{\int_{\varphi}^{\Sigma}Y}^{\tau,i}
      &\quad\text{défini par la matrice }M(\tau,i,Y).
  \end{aligned}
  \tag{2.2.3.3}
\end{equation}

\VerdierSemanticAnchor{ch1-2-2-3-pi-extension}
Le foncteur $\int_{\varphi}^{\Pi}$ est appelé la
$\Pi$-\emph{extension} du foncteur $\int_\varphi$.

\VerdierSemanticAnchor{ch1-2-2-3-sigma-extension}
Le foncteur $\int_{\varphi}^{\Sigma}$ est appelé la
$\Sigma$-\emph{extension} du foncteur $\int_\varphi$. Ces deux extensions
seront utilisées par la suite.

\VerdierSemanticAnchor{ch1-2-2-3-natural-map}
On a, d'ailleurs, un morphisme naturel de foncteurs
$\int_{\varphi}^{\Sigma}\longrightarrow\int_{\varphi}^{\Pi}$
correspondant au morphisme fonctoriel de la somme directe dans le produit.

\VerdierSemanticAnchor{ch1-2-2-4}
\noindent\textbf{2.2.4.}\quad
Soient
$[n]\xrightarrow{\varphi}[m]\xrightarrow{\psi}[r]$ deux applications et
supposons, pour fixer les idées, que les produits dénombrables soient
représentables dans $C$. On a alors les relations :

\VerdierSemanticAnchor{ch1-2-2-4-eq1}
\begin{equation}
  \int_{\psi}^{\Pi}\int_{\varphi}^{\Pi}
  =\int_{\psi\varphi}^{\Pi},
  \tag{2.2.4.1}
\end{equation}

\VerdierSemanticAnchor{ch1-2-2-4-eq2}
\begin{equation}
  \int_{\operatorname{id}}^{\Pi}=\operatorname{id}.
  \tag{2.2.4.2}
\end{equation}

\VerdierSemanticAnchor{ch1-2-2-4-editorial-caveat}
Le rédacteur se voit ici dans la pénible obligation de tirer l'irritant coup
de chapeau habituel. Le signe $=$ de la formule (2.2.4.1) n'est valable que
lorsqu'on dispose d'un foncteur produit possédant des propriétés
d'associativité convenables, ce qui n'est jamais le cas dans la pratique.

\VerdierSemanticAnchor{ch1-2-2-4-canonical-isomorphism}
Dans le cas général, le signe ``$=$'' doit être remplacé par un isomorphisme
canonique qui vérifie les propriétés de compatibilité usuelles du type
(1.2.1.3).

\VerdierSemanticAnchor{ch1-2-2-4-sigma-analogue}
Des relations analogues aux relations (2.2.4.1) sont valables pour le
foncteur $\int_{\varphi}^{\Sigma}$.

\VerdierSemanticAnchor{ch1-2-2-4-summable-restriction}
Lorsque les produits et sommes dénombrables ne sont pas nécessairement
représentables, la formule (2.2.4.1) est encore valable lorsqu'on se
restreint aux complexes $\psi\varphi$-sommables.

\VerdierSemanticAnchor{ch1-2-2-4-simple-complex}
Lorsque $\psi$ est l'unique morphisme $[n]\longrightarrow[1]$, on utilise
la notation $\int_n$. Pour tout complexe $n$-uple sommable $Y$, le complexe
$\int_n Y$ est appelé le \emph{complexe simple associé au complexe $Y$}.
